\documentclass[10pt,a4paper]{amsart}
\usepackage[margin=19mm]{geometry}
\usepackage{amsmath,amssymb,mathtools}
\usepackage[scr=rsfs,cal=euler]{mathalfa}
\usepackage{xfrac}
\usepackage[unicode,hidelinks,bookmarks=false]{hyperref}
\newcommand{\ie}{i.e.\ }
\newcommand{\cf}{cf.\ }
\newcommand{\resp}{respectively\ }
\newcommand{\Subsection}{\subsection}
\providecommand{\nobreakdash}{\nobreak\hbox{-}}
\setlength{\emergencystretch}{2em}
\multlinegap0pt
\usepackage{xcolor}
\usepackage{amssymb}
\usepackage{mathtools}
\usepackage{bm}
\usepackage[shortlabels]{enumitem}
\newtheorem{theorem}{Theorem}
\newcommand{\Hd}{\mathbb{H}_d}
\newcommand{\Hm}{\mathbb{H}_m}
\newcommand{\inner}[2]{\langle #1, #2 \rangle}
\newcommand{\op}{\mathrm{op}}
\title{Symplectic Neural Operators for Learning Infinite Dimensional Hamiltonian Systems}
\author{Yeang Makara, Yusuke Tanaka, Takashi Matsubara, Takaharu Yaguchi}
\date{Source: arXiv:2605.15881v1 (CC BY 4.0)}
\begin{document}
\maketitle

\noindent\emph{Source and licence.} Symplectic Neural Operators for Learning Infinite Dimensional Hamiltonian Systems. Version arXiv:2605.15881v1 (CC BY 4.0), distributed by arXiv under the Creative Commons Attribution 4.0 International licence, http://creativecommons.org/licenses/by/4.0/, as recorded in the arXiv licence metadata for this exact version and shown on the abstract page https://arxiv.org/abs/2605.15881v1. Full licence notice: \texttt{licenses/arxiv-2605.15881-CC-BY-4.0.txt}.\par\smallskip

\noindent\textit{Editorial scope.} Counted unit: the article's theorem nsafno (source0.tex lines 802--816). The article's complete original proof of it is delivered with it (source0.tex lines 1355--1489, one \texttt{proof} environment of 135 lines, which the article places away from the statement). No mathematics is written, repaired or re-derived: the statement and the proof are the article's own lines, in the article's own notation, and no retained line is substituted or re-flowed.  No further source-local context is needed: every cross-reference of every retained line resolves inside the two delivered spans.  Nothing was omitted from any retained span, so the package carries no omission record.  No retained line contains a citation, so the wrapper carries no bibliography and no bibliography entry.  No retained line contains a figure environment, an image-inclusion command or drawing code, so no image is delivered and the package declares no asset dependency.  Editorial formatting only, in addition: an AMS \texttt{amsart} preamble that restates the article's own declarations and macros used by the retained lines, namely the environments \texttt{theorem} on separate counters, together with the standard packages those lines need.  The retained environments print in this document as Theorem 1 in order of appearance, whereas the article prints the same items with the numbering of its own full document.  The complete native source of the article as distributed in the arXiv e-print for this exact version is bundled unchanged as \texttt{sources/200688/source0.tex}.
\begin{theorem}[Gradient formula and self-adjoint linearization]\label{nsafno}
Let $\Hd=L^2(\mathbb{T}^n;\mathbb{R}^d)$ and $\Hm=L^2(\mathbb{T}^n;\mathbb{R}^m)$, and let $K:\Hd\to\Hm$ be a bounded linear operator with adjoint $K^*$. Let $\Phi_\theta:\mathbb{R}^m\to\mathbb{R}$ be $C^2$ with $\sup_{z}\|D^2\Phi_\theta(z)\|_{\mathrm{op}}<\infty$. Define energy functional 
\(
\mathcal{E}_\theta(u)=\int_{\mathbb{T}^n}\Phi_\theta\big((Ku)(x)\big)\,dx.
\)
Then $\mathcal{E}_\theta$ is Fr\'echet differentiable with
\(
\nabla\mathcal{E}_\theta(u)=K^*\big(\nabla\Phi_\theta(Ku)\big),
\)
and $\mathcal{G}_\theta:=\nabla\mathcal{E}_\theta$ is Fr\'echet differentiable with
\(
D\mathcal{G}_\theta(u)h = K^*\!\left(D^2\Phi_\theta(Ku)\,(Kh)\right), h\in\Hd,
\)
equivalently $D\mathcal{G}_\theta(u)=K^* M_{H_\theta(u)} K$, where $H_\theta(u)(x)=D^2\Phi_\theta((Ku)(x))$; in particular, $D\mathcal{G}_\theta(u)$ is a bounded self-adjoint operator on $\Hd$ for every $u\in\Hd$.
\end{theorem}

\begin{proof}(Theorem \ref{nsafno})
Let $\Omega=\mathbb{T}^n$.

Step 1 (A pointwise Lipschitz bound).

Since $\sup_{z}\|D^2\Phi_\theta(z)\|_{\mathrm{op}}\le M$, the gradient is globally Lipschitz:
for all $z_1,z_2\in\mathbb{R}^m$,
\[
\|\nabla\Phi_\theta(z_1)-\nabla\Phi_\theta(z_2)\|
\le M\|z_1-z_2\|.
\]
Also, by the fundamental theorem of calculus along the segment $t\mapsto tz$,
\[
\nabla\Phi_\theta(z)-\nabla\Phi_\theta(0)
=\int_0^1 D^2\Phi_\theta(tz)\,z\,dt,
\]
hence
\[
\|\nabla\Phi_\theta(z)\|\le \|\nabla\Phi_\theta(0)\| + M\|z\|.
\]
Therefore, for any $v\in \Hm$ we have $\nabla\Phi_\theta(v)\in \Hm$ and
\[
\|\nabla\Phi_\theta(v)\|_{\Hm}
\le \|\nabla\Phi_\theta(0)\|\,|\Omega|^{1/2} + M\|v\|_{\Hm},
\]
where $|\Omega|=1$ for Haar probability measure on $\mathbb{T}^n$.

\textbf{Step 2 (Fr\'echet differentiability of $\mathcal{E}_\theta$ and gradient formula).}
Fix $u,h\in\Hd$ and set $v:=Ku\in\Hm$ and $w:=Kh\in\Hm$.
For each $x\in\Omega$, define $\psi_x(t):=\Phi_\theta(v(x)+t w(x))$.
By the one-dimensional fundamental theorem of calculus,
\[
\Phi_\theta(v+w)-\Phi_\theta(v)
=\int_0^1 \nabla\Phi_\theta(v+t w)\cdot w\,dt,
\]
where ``$\cdot$'' is the Euclidean dot product in $\mathbb{R}^m$.
Hence
\[
\mathcal{E}_\theta(u+h)-\mathcal{E}_\theta(u)
=\int_\Omega\int_0^1 \nabla\Phi_\theta(v+t w)\cdot w\,dt\,dx.
\]
Subtract the candidate derivative $\int_\Omega \nabla\Phi_\theta(v)\cdot w\,dx$:
\[
R(u,h)
:=\mathcal{E}_\theta(u+h)-\mathcal{E}_\theta(u)-\int_\Omega \nabla\Phi_\theta(v)\cdot w\,dx
=\int_\Omega\int_0^1\big(\nabla\Phi_\theta(v+t w)-\nabla\Phi_\theta(v)\big)\cdot w\,dt\,dx.
\]
Using the Lipschitz bound from Step 1,
\[
\big|\big(\nabla\Phi_\theta(v+t w)-\nabla\Phi_\theta(v)\big)\cdot w\big|
\le \|\nabla\Phi_\theta(v+t w)-\nabla\Phi_\theta(v)\|\;\|w\|
\le M t \|w\|^2.
\]
Integrating in $t$ and $x$ gives
\[
|R(u,h)|
\le \int_\Omega\int_0^1 M t \|w(x)\|^2\,dt\,dx
=\frac{M}{2}\|w\|_{\Hm}^2
\le \frac{M}{2}\|K\|^2\|h\|_{\Hd}^2.
\]
Therefore $R(u,h)=o(\|h\|_{\Hd})$ as $h\to 0$, and $\mathcal{E}_\theta$ is Fr\'echet differentiable with
\[
D\mathcal{E}_\theta(u)[h]=\int_\Omega \nabla\Phi_\theta((Ku)(x))\cdot (Kh)(x)\,dx
=\inner{\nabla\Phi_\theta(Ku)}{Kh}_{\Hm}.
\]
By definition of the adjoint, $\inner{\nabla\Phi_\theta(Ku)}{Kh}_{\Hm}
=\inner{K^*\nabla\Phi_\theta(Ku)}{h}_{\Hd}$, hence
\[
\nabla\mathcal{E}_\theta(u)=K^*\big(\nabla\Phi_\theta(Ku)\big).
\]

Step 3 (Fr\'echet differentiability of $\mathcal{G}_\theta$ and formula for $D\mathcal{G}_\theta(u)$.

We have $\mathcal{G}_\theta(u)=K^*F(Ku)$ where $F(v):=\nabla\Phi_\theta(v)$ acts pointwise.
Since $K^*$ and $K$ are bounded linear, it suffices to show $F:\Hm\to\Hm$ is Fr\'echet differentiable and compute $DF(v)$.

Fix $v\in\Hm$ and $w\in\Hm$. Pointwise in $\mathbb{R}^m$,
\[
\nabla\Phi_\theta(v+w)-\nabla\Phi_\theta(v) - D^2\Phi_\theta(v)\,w
= \int_0^1 \big(D^2\Phi_\theta(v+t w)-D^2\Phi_\theta(v)\big)\,w\,dt.
\]
Assuming $\Phi_\theta\in C^2$, the map $z\mapsto D^2\Phi_\theta(z)$ is continuous.
Moreover, by boundedness of $D^2\Phi_\theta$,
\[
\|\big(D^2\Phi_\theta(v+t w)-D^2\Phi_\theta(v)\big)\,w\|
\le 2M\|w\|.
\]
Using dominated convergence (in $x$ and $t$) and $w\in L^2$, one obtains
\[
\frac{\|F(v+w)-F(v)-M_{D^2\Phi_\theta(v)}w\|_{\Hm}}{\|w\|_{\Hm}}\to 0
\quad\text{as }\|w\|_{\Hm}\to 0,
\]
so $F$ is Fr\'echet differentiable with
\[
(DF(v)w)(x)=D^2\Phi_\theta(v(x))\,w(x).
\]
Applying the chain rule,
\[
D\mathcal{G}_\theta(u)h
=K^*\big(DF(Ku)(Kh)\big)
=K^*\big(D^2\Phi_\theta(Ku)\,(Kh)\big),
\]
which is the claimed formula.

Step 4 (Boundedness and self-adjointness).

Define the multiplication operator on $\Hm$:
\[
(M_{H_\theta(u)}\varphi)(x):=H_\theta(u)(x)\,\varphi(x),
\quad H_\theta(u)(x)=D^2\Phi_\theta((Ku)(x)).
\]
Since $\|H_\theta(u)(x)\|_{\op}\le M$ for all $x$, we have
\[
\|M_{H_\theta(u)}\varphi\|_{\Hm}\le M\|\varphi\|_{\Hm},
\]
so $M_{H_\theta(u)}$ is bounded.
Thus $D\mathcal{G}_\theta(u)=K^*M_{H_\theta(u)}K$ is bounded.

Finally, for $h_1,h_2\in\Hd$,
\[
\inner{D\mathcal{G}_\theta(u)h_1}{h_2}_{\Hd}
=\inner{M_{H_\theta(u)}Kh_1}{Kh_2}_{\Hm}
=\int_\Omega (Kh_1(x))^\top H_\theta(u)(x)\,(Kh_2(x))\,dx.
\]
Since $H_\theta(u)(x)$ is symmetric for every $x$, the integrand satisfies
\[
(Kh_1)^\top H_\theta(u)\,(Kh_2)=(Kh_2)^\top H_\theta(u)\,(Kh_1),
\]
hence
\[
\inner{D\mathcal{G}_\theta(u)h_1}{h_2}_{\Hd}
=\inner{h_1}{D\mathcal{G}_\theta(u)h_2}_{\Hd}.
\]
Therefore $D\mathcal{G}_\theta(u)$ is self-adjoint.
\end{proof}

\end{document}
